\section{Results}

\subsection{Functional Outcomes}
Table~\ref{tab:results} summarizes the main outcomes observed in the final validated configuration.

\begin{table}[h]
\centering
\caption{End-to-end validation outcomes.}
\label{tab:results}
\begin{tabular}{p{0.47\linewidth} p{0.14\linewidth} p{0.30\linewidth}}
\toprule
\textbf{Checkpoint} & \textbf{Result} & \textbf{Evidence} \\
\midrule
Platform services reachable in stable experimental mode & Pass & API server and gateway pods remained \texttt{Running}; measurement prerequisites re-established before each run \\
TLS enrollment handshake completion & Pass & Repeated enrollment request/ack/completion sequences across all trials \\
Gateway-visible heartbeat supervision & Pass & Fresh gateway-runtime heartbeat timestamps and positive connected state in every trial \\
Gateway-mediated application activation & Pass & Application CRD reached \texttt{Running} phase with authoritative controller confirmation for the target device in every trial \\
Transactional trial completion & Pass & Enrollment, heartbeat supervision, and deployment all succeeded in the same trial for all 50 trials \\
Measurement artifact reproducibility & Pass & Machine-readable campaign summary and generated figures reproduced from the stored JSON artifact \\
\bottomrule
\end{tabular}
\end{table}

The table provides a compact view of the experimental outcome, but its significance lies in the fact that the checkpoints cover the entire orchestration path rather than a single protocol stage. The platform does not merely accept a connection; it preserves the semantic coherence of enrollment, liveness supervision, and device-side application activation across multiple components and across repeated trials.

\subsection{Measured Campaign Results}
The final measurement campaign executed 50 independent trials for enrollment, heartbeat supervision, and deployment, plus one system snapshot run. Each trial was then composed into a transactional record that is successful only when all three stages succeed. The resulting summary is reported in Tables~\ref{tab:quantitative-success} and~\ref{tab:quantitative-latency}, and visualized in Figures~\ref{fig:latency-ci-profile} and~\ref{fig:latency-boxplot}.

All three measured stages achieved a success rate of 1.0 over 50 trials, with the same result for the transactional metric. Wilson score intervals are bounded by a lower limit of 92.9\%, improving upon the 88.6\% attainable with $n=30$. Mean enrollment latency was 539.5~ms, mean heartbeat observability latency was 3.90~ms, mean deployment latency was 521.0~ms, and mean end-to-end latency was 1060.5~ms. The corresponding 95\% confidence interval for end-to-end latency was 1022.8--1098.2~ms, computed with $t_{0.975,\,49}=2.010$.

Deployment success is assessed directly from the Application CRD \texttt{phase} field reaching \texttt{Running}, which is now the authoritative signal because the application controller is active and its reconciliation loop has been hardened. The API server statistics are retained as a secondary corroborating source and are preserved in the raw record for reproducibility.

\begin{table}[h]
\centering
\caption{Success-oriented summary from the final 50-trial campaign.}
\label{tab:quantitative-success}
\begin{tabular}{p{0.28\linewidth} p{0.09\linewidth} p{0.30\linewidth} p{0.16\linewidth}}
\toprule
\textbf{Metric} & \textbf{n} & \textbf{Success rate} & \textbf{Goodput} \\
\midrule
Enrollment & 50 & 100\% [92.9\%, 100\%] & 1.85 s$^{-1}$ \\
Heartbeat supervision & 50 & 100\% [92.9\%, 100\%] & 256.47 s$^{-1}$ \\
Deployment & 50 & 100\% [92.9\%, 100\%] & 1.92 s$^{-1}$ \\
Transactional completion & 50 & 100\% [92.9\%, 100\%] & -- \\
System snapshot & 1 & 100\% [20.7\%, 100\%] & -- \\
\bottomrule
\end{tabular}
\end{table}

Because all stage-level success rates are equal to 1.0 over the same sample size ($n=50$), Wilson confidence intervals become numerically identical across stages. For this reason, success rates are kept as a compact table rather than as the primary figure, and the graphical analysis focuses on latency distributions and uncertainty, which are more discriminative in the present campaign.

\begin{table}[h]
\centering
\caption{Latency profile from the final 50-trial campaign ($t_{0.975,\,49}=2.010$).}
\label{tab:quantitative-latency}
\begin{tabular}{p{0.22\linewidth} p{0.12\linewidth} p{0.12\linewidth} p{0.12\linewidth} p{0.12\linewidth} p{0.10\linewidth} p{0.17\linewidth}}
\toprule
\textbf{Stage} & \textbf{Mean} & \textbf{Median} & \textbf{p95} & \textbf{p99} & \textbf{CV} & \textbf{95\% CI} \\
\midrule
Enrollment & 539.5 ms & 530.9 ms & 686.9 ms & 721.4 ms & 0.135 & [518.8, 560.2] ms \\
Heartbeat & 3.90 ms & 3.73 ms & 4.90 ms & 6.82 ms & 0.210 & [3.67, 4.13] ms \\
Deployment & 521.0 ms & 477.4 ms & 840.6 ms & 880.5 ms & 0.230 & [486.9, 555.1] ms \\
End-to-end & 1060.5 ms & 1019.2 ms & 1370.1 ms & 1424.2 ms & 0.125 & [1022.8, 1098.2] ms \\
\bottomrule
\end{tabular}
\end{table}

\begin{figure}[h]
\centering
\includegraphics[width=\columnwidth]{figures/latency_ci_profile.png}
\caption{Latency profile by stage ($n=50$). Bars show sample means, error bars show 95\% confidence intervals of the mean, and triangle markers denote p95.}
\label{fig:latency-ci-profile}
\end{figure}

\begin{figure}[h]
\centering
\includegraphics[width=\columnwidth]{figures/latency_boxplot.png}
\caption{Distribution of trial-level latencies on a logarithmic scale ($n=50$). Diamond markers indicate sample means.}
\label{fig:latency-boxplot}
\end{figure}

The boxplot reveals that heartbeat observability is two orders of magnitude faster than enrollment or deployment, as expected, and that the end-to-end latency remains tightly concentrated around 1.06~s despite visible tail growth in the deployment component. The deployment tail dominates the transactional tail, which suggests that future optimization effort should focus on WASM-module activation orchestration rather than on enrollment or liveness supervision. The coefficient of variation for enrollment (CV\,=\,0.135) is the lowest among non-trivial stages, confirming that the TLS and CBOR handshake path is stable. Deployment shows the highest CV\,=\,0.230, reflecting the asymmetric tail introduced by WASM instantiation on the embedded target.

\subsection{How to Read the Metrics}
The reported means should be interpreted together with confidence intervals and variability indicators rather than in isolation. A low average is useful only if the interval remains narrow, the percentile tail remains controlled, and the success rate stays close to one. For that reason, the manuscript treats latency and correctness as jointly relevant outcomes and introduces the transactional success metric to summarize whether the full path from secure attachment to application activation completed correctly.

\subsection{Observed Session Behavior}
During active sessions, the device remained connected from the gateway's authoritative runtime perspective while heartbeats were periodically transmitted and observed. The measured heartbeat latency should therefore be read as liveness-observation latency, that is, the time needed by the northbound harness to verify the freshness of the most recent heartbeat, rather than as the heartbeat emission period itself. This distinction is important because the system goal is not to minimize the firmware heartbeat interval, but to guarantee that liveness remains quickly and reliably observable to the control plane.

Across the 50-trial campaign, no enrollment, heartbeat-supervision, or deployment failures were observed. The corresponding percentile profiles remained tight, indicating that the platform is not merely correct on average but also stable across repetitions. The enrollment component exhibited the smallest tail spread among the nontrivial workflow stages; the deployment tail is wider but controlled, with p99 at 880.5~ms.

\subsection{Impact of Reconciliation Hardening}
One of the most relevant findings of the study is that the dominant failure modes were neither cryptographic nor transport-level. Instead, they emerged from interactions between controller logic, status-field naming, and control-plane race conditions. The final quantitative campaign was only possible after hardening the reconcile loop, correcting the CRD snake\_case status field mismatch, and ensuring the application controller correctly propagates the authoritative phase back to the cluster when in-memory state diverges from persisted state.

This outcome strengthens the paper's systems claim. It shows that operational correctness in cyber-physical orchestration cannot be reduced to protocol success alone. Control-plane semantics matter, and they can invalidate an otherwise correct protocol path if reconciliation is overly aggressive, if status ownership is ambiguous, or if field naming conventions are mismatched between writers.

\subsection{Practical Interpretation}
From an engineering perspective, the results indicate that RETROSPECT can support a realistic orchestration workflow for embedded experimentation and early-stage platform deployment. The combination of declarative cloud control, secure gateway-mediated attachment, reproducible validation on a commodity K3s environment, and transaction-oriented metrics makes the platform useful both as a research artifact and as a baseline for future industrial prototyping.

The present results should be interpreted as evidence of functional maturity together with first-order performance characterization. They demonstrate that the architecture works as intended under the targeted workflow, that the main orchestration semantics are stable enough to support subsequent experimentation on scale and hardware diversity, and that the confidence intervals are tight enough to detect meaningful latency regressions in future work.
